% chapters/01-intro.tex -- Вступ (ДСТУ 3008:2015, п. 5.1; Рекомендації
% КПІ уточнюють для магістерської дисертації, що обсяг вступу не має
% перевищувати 5% основної частини).
\kpithesisUnnumberedChapter{Вступ}

\textbf{Актуальність теми.} \dots

\textbf{Мета і завдання дослідження.} \dots
Для досягнення мети потрібно розв'язати такі завдання:
\begin{itemize}
  \item \dots;
  \item \dots;
  \item \dots.
\end{itemize}

\textbf{Об'єкт дослідження.} \dots

\textbf{Предмет дослідження.} \dots

\textbf{Методи дослідження.} \dots

\textbf{Наукова новизна одержаних результатів} (обов'язково для
освітньо-наукової програми). \dots

\textbf{Практичне значення одержаних результатів.} \dots

\textbf{Апробація результатів.} \dots

\textbf{Публікації.} \dots
